\scanpage{76}
(L'indétermination est donc section d'un rev. principal de $S$ de groupe $\mu_2$,
défini mod él. de $\mu_2(S)\ \ldots$).
Donc $s_0$ est alors donné par la formule $(*)$.
(Rien n'empêche d'ailleurs que $f_V(s_0)=0$)

\deleted{(E) \unsure{Supposons à présent}}
Si on se donne une hom
\[
\varphi_G^C:G\mathbin{-}\operatorname{Aut}(C,Q_C)
\]
[$C$ fibré proj. de rg $n$, $Q_C$ \inserted{hyper}quadrique]
satisfaisant les conditions

$1)$ Les $\tau_i^X$ réflexions \deleted{orth}
\hfill$(0\leq i\leq n)$

\hfill i.e. $h_i\in\Gamma C$, $H_i$ hyperplans)

$2)$ Les $h_i$ engendrent $C$, les
$\bigcap_{0\leq i\leq n}H_i^{\unsure{C}}=\varnothing$

$3)$ \deleted{\unsure{H_i}}
\quad$h_i\notin H_{i+1}$\hfill$0\leq i\leq n-1$

\deleted{et $4)$ \unsure{un}
$\unsure{L_0|h_0}\notin Q_C$ en chaque fibre}

\margin{Pour $n\geq2$}
\margin{(Si $n=1$, une réflexion de $C$ donne \unsure{pour}
$C=P(W)$, \unsure{ne remonte pas} de \uline{façon unique}
en \unsure{réfl.} de $V\ldots$]}

alors loc. (et) ça provient de
$(V,f_V,(\tau_i^V))$ satisfaisant les conditions de (D),
$V$ déterm. mod $\otimes L$ ($L$ inv),
$f_V$ défini mod unité.
On utilise \unsure{$h_0$ pour rigidifier};
$L_0\subset V$, on les ``épingle'' i.e éliminer l'indét. $\otimes L$,
en \unsure{plongeant|choisissant} $L\simeq\mathcal O_S$
i.e. section $a_0$ de $L_0$.

\mspage{75}

